\section*{附加习题: 局部欧氏空间}
\addcontentsline{toc}{section}{附加习题: 局部欧氏空间}
\label{section:ch08-extra}

\begin{remark}
一个空间 $X$ 称为\textbf{局部 $m$-维欧氏的}(locally $m$-Euclidean), 如果对于每一个 $x\in X$, 存在 $x$ 的一个邻域同胚于 $\mathbb{R}^{m}$ 中的一个开集. 这样的空间自然满足 $T_1$ 公理, 但不一定是 Hausdorff 空间. (见第 36 节的习题.) 若 $X$ 还是一个 Hausdorff 空间, 并且有可数基, 那么称 $X$ 为一个 $m$-\textbf{维流形}($m$-manifold).

以下习题中统统假设 $X$ 是局部 $m$-维欧氏空间.
\end{remark}

\begin{exercise}
证明 $X$ 是局部紧致的和局部可度量化的.
\end{exercise}

\begin{exercise}
考虑关于 $X$ 的以下条件:
\begin{enumerate}
\item $X$ 是一个紧致的 Hausdorff 空间.
\item $X$ 是一个 $m$-维流形.
\item $X$ 是可度量化的.
\item $X$ 是正规的.
\item $X$ 是 Hausdorff 的.
\end{enumerate}
证明: (i)$\Rightarrow$(ii)$\Rightarrow$(iii)$\Rightarrow$(iv)$\Rightarrow$(v).
\end{exercise}

\begin{exercise}
证明 $\mathbb{R}$ 是局部 $1$-维欧氏空间, 并满足 (ii), 但不满足 (i).
\end{exercise}

\begin{exercise}
证明 $\mathbb{R}\times\mathbb{R}$ 在字典序拓扑下是局部 $1$-维欧氏空间, 并且满足 (iii), 但不满足 (ii).
\end{exercise}

\begin{exercise}
证明长直线是局部 $1$-维欧氏空间, 并且满足 (iv), 但不满足 (iii). (见第 24 节的习题.)
\end{exercise}

\begin{exercise}
存在一个空间, 它是局部 $2$-维欧氏空间, 满足 (v), 但不满足 (iv). 这个空间的构造方法如下: 设 $A$ 是 $\mathbb{R}^{3}$ 的以下子空间:
\[
A=\{\,(x,y,0)\mid x>0\,\}.
\]
给定实数 $c$, 令 $B_c$ 是 $\mathbb{R}^{3}$ 的以下子空间:
\[
B_{c}=\{\,(x,y,c)\mid x\leqslant 0\,\}.
\]
设 $X$ 为 $A$ 和所有子空间 $B_c$ 的并, 其中 $c$ 取遍所有实数. 选取以下三种集合作为基将 $X$ 拓扑化:
\begin{enumerate}
\item $U$, 其中 $U$ 在 $A$ 中是开的.
\item $V$, 其中 $V$ 是子空间 $B_c$ 中所有满足 $x<0$ 的点构成的 $B_c$ 中的一个开集.
\item 对于 $\mathbb{R}$ 的每一个开区间 $I=(a,b)$、每一个实数 $c$ 以及每一个 $\varepsilon>0$, 集合 $A_c(I,\varepsilon)\cup B_c(I,\varepsilon)$, 其中
\[
A_{c}(I,\varepsilon)=\{\,(x,y,0)\mid 0<x<\varepsilon\text{ 和 }c+ax<y<c+bx\,\},
\]
\[
B_{c}(I,\varepsilon)=\{\,(x,y,c)\mid -\varepsilon<x\leqslant 0\text{ 和 }a<y<b\,\}.
\]
\end{enumerate}
空间 $X$ 称为"Pr\"ufer 流形".
\begin{enumerate}
\item 画出 $A_c(I,\varepsilon)$ 和 $B_c(I,\varepsilon)$ 的草图.
\item 证明 (i)~(iii) 三种集合共同构成了 $X$ 的拓扑的一个基.
\item 证明: 定义为
\[
f_{c}(x,y)=
\begin{cases}
(x,c+xy,0) & \text{对于 }x>0,\\
(x,y,c) & \text{对于 }x\leqslant 0
\end{cases}
\]
的映射 $f_c\colon\mathbb{R}^{2}\to X$ 是 $\mathbb{R}^{2}$ 和 $X$ 的子空间 $A\cup B_c$ 之间的一个同胚.
\item 证明 $A\cup B_c$ 在 $X$ 中是开的. 证明 $X$ 是 $2$-维欧氏空间.
\item 证明 $X$ 是 Hausdorff 的.
\item 证明 $X$ 不是正规的. [提示: $X$ 的子空间
\[
L=\{\,(0,0,c)\mid c\in\mathbb{R}\,\}
\]
是闭的和离散的. 与第 31 节中的习题 3 比较.]
\end{enumerate}
\end{exercise}

\begin{exercise}
证明: $X$ 是 Hausdorff 的当且仅当 $X$ 是完全正则的.
\end{exercise}

\begin{exercise}
证明: $X$ 是可度量化的当且仅当 $X$ 是仿紧致的和 Hausdorff 的.
\end{exercise}

\begin{exercise}
证明: 若 $X$ 是可度量化的, 则 $X$ 的每一个分支都是一个 $m$-维流形.
\end{exercise}
